\documentclass[10pt]{letter}
\usepackage[T1]{fontenc}
\usepackage[margin=0.68in]{geometry}
\setlength{\emergencystretch}{3em}
\setlength{\longindentation}{0.48\textwidth}
\signature{Lizhen Shao\\Corresponding author\\University of Science and Technology Beijing\\\texttt{lshao@ustb.edu.cn}}
\address{Department of Control Science and Engineering\\University of Science and Technology Beijing\\No.\ 30 Xueyuan Road, Haidian District\\Beijing 100083, China}
\date{\today}

\begin{document}

\begin{letter}{Editor-in-Chief\\Computers \& Industrial Engineering}

\opening{Dear Editor-in-Chief,}

Please consider our Research Paper entitled ``Two-stage MILP scheduling and
structure-aware learning-enhanced branch-and-cut for dual-source mixed-flow
rotating-chamber cluster tools'' for publication in \emph{Computers \&
Industrial Engineering}.

The paper addresses an industrial scheduling problem arising in rotating-
chamber cluster tools used in a PECVD process setting. Product wafers and
reusable vacuum-zone dummy wafers share chambers, load locks, and robots while
obeying recipe compatibility, residency, cleaning, capacity, and transfer
constraints. Their concurrent interaction creates a finite-horizon scheduling
problem that is not captured by single-flow cluster-tool formulations.

The paper makes two connected contributions. First, it develops a finite-
horizon two-stage MILP that integrates canonical same-recipe pairing and slot
incidence with transfers, cleaning epochs, reusable dummy-wafer allocation,
shared-resource ordering, residence limits, and event timing. The canonical
incidence removes equivalent labels; the second stage preserves the stage-one
discrete schedule and makespan within numerical tolerance while refining
physical route waits and processing cadence. Second, it develops RDMCT-HBS, a
structure-aware learning-enhanced branch-and-cut method combining a
precomputed SPBS incumbent, reconstructed variable-role features, A3C-inspired
parallel policy-gradient training, bounded beam decoding, and role-submodular
completion.

The computational evidence is both extensive and mechanism-specific. In a
2700-run nominal comparison, RDMCT-HBS records the lowest observed mean
primal--dual integral and mean solver time among seven methods. A separate
960-run paired experiment shows that supplying a precomputed SPBS incumbent
reduces mean primal--dual integral by 26.19\% and increases incumbent
availability from 51.67\% to 100\% during the timed SCIP stage; offline SPBS
construction is excluded from that budget. In 1500 matched runs under the same
600-s total budget, the complete two-stage workflow returns schedules with
physical route-total wait, route-maximum wait, cadence coefficient of
variation, and cadence deviation lower by 51.91\%, 68.89\%, 39.70\%, and
94.05\%, respectively, than the no-refinement workflow. A complete 300-run
factorial experiment further shows that wafer count and mixed-flow recipe
composition dominate the observed difficulty variation. The
larger-scale stress test adds 729 time-limited runs on nine 40--64-wafer
configurations, every one of which returns an independently validated
incumbent. Across the complete evidence programme, all 6273 planned runs are
accounted for and all 6041 saved incumbents pass independent feasibility
checking against their frozen generated MILPs.

The work fits the journal because it integrates a new industrial-engineering
model with an explainable, exact-optimization-compatible learning method and
evaluates both computational and manufacturing consequences. The formulation
and solution architecture are also transferable to other tightly coupled
multi-resource scheduling systems with auxiliary flows and shared transport.

% Before submission, the authors must confirm originality, exclusive
% consideration, approval by all authors, funding, acknowledgements, and
% individual CRediT roles. These facts are intentionally not asserted here
% until author confirmation is received.

Thank you for your consideration.

\closing{Sincerely,}

\end{letter}
\end{document}
